\chapter{Compactness}

\section{Uniform Continuity}

\begin{definition}{Uniform continuity}
    $f: (M, d) \rightarrow (N, \rho)$ is uniformly continuous if for every $\epsilon > 0$, there exists $\delta > 0$ such that $\rho(f(x), f(y)) < \epsilon$ whenever $d(x, y) < \delta$. 
\end{definition}

\begin{exercise}{8.44}
    Any Lipschitz map is in fact uniformly continuous.
    This implies that isometries are also uniformly continuous.
\end{exercise}

\begin{exercise}{8.45}
    \textit{Every} map $f: \mathbb{N} \rightarrow \mathbb{R}$ is uniformly continuous.
    Intuitively, this is because natural numbers never get ``arbitrarily close'' to each other.
\end{exercise}

\begin{exercise}{48}
    A uniformly continuous map sends Cauchy sequences to Cauchy sequences.
\end{exercise}

\begin{exercise}{8.54}
    In bounded, noncompact subsets of $\mathbb{R}$, continuous functions need not uniformly continuous.
\end{exercise}

\begin{exercise}{8.55}
    Bounded continuous maps in $\mathbb{R}$ need not be uniformly continuous: for example, $f(x) = \sin(e^x)$, which oscillates ``increasingly fast''.
    Furthermore, unbounded continuous functions \textit{can} be uniformly continuous, as easily seen by the identity function.
\end{exercise}

\begin{exercise}{8.56}
    \textbf{A criterion for uniform continuity}: $f: (M, d) \rightarrow (N, \rho)$ is uniformly continuous iff $\rho(f(x_n), f(y_n)) \rightarrow 0$ for any pair of sequences $(x_n), (y_n) \subset M$ such that $d(x_n, y_n) \rightarrow 0$.
\end{exercise}

\begin{exercise}{8.57}
    A type of Lipschitz continuity: If $\lvert f(x) - f(y) \rvert \leq K \lvert x - y \rvert^a$ for $a > 0$ then $f$ is uniformly continuous.
\end{exercise}

\begin{exercise}{8.58}
    Any function $f: \mathbb{R} \rightarrow \mathbb{R}$ with a bounded derivative is Lipschitz of order 1.
\end{exercise}

\begin{exercise}{8.59}
    A function satisfying a Lipschitz condition as in Ex. 8.57 for $a > 1$ is \textbf{necessarily constant}.
\end{exercise}

\begin{exercise}{8.61}
    In general, we know that completeness is not necessarily maintained by homeomorphisms.
    However, if $f, f^{-1}$ are \textit{uniformly continuous} homeomorphisms (\textbf{uniform homeomorphisms}) between two metric spaces, then completeness is preserved.
    We call such metric spaces $\textbf{uniformly homeomorphic}$.
\end{exercise}

\begin{proposition}{8.14}
    If $f: M \rightarrow N$ is uniformly continuous, then $f$ maps totally bounded sets into totally bounded sets.
\end{proposition}

\begin{theorem}{8.15}
    If $M$ is a compact metric space, then every continuous map $f: M \rightarrow N$ is uniformly continuous.
\end{theorem}

\begin{theorem}{8.16}
    Let $D$ be dense in $M$, let $N$ be complete, and let $f: D \rightarrow N$ be uniformly continuous.
    Then, $f$ extends uniquely to a uniformly continuous map $F: M \rightarrow N$, defined on all of $M$.
    Moreover, if $f$ is an isometry, then so is the extension $F$.
\end{theorem}

\begin{corollary}{8.17}
    Completions are unique (up to isometry).
    That is, if $M_1, M_2$ are completions of $M$, then $M_1, M_2$ are isometric.
\end{corollary}

\section{Equivalent Metrics}

\begin{proposition}{8.19}
    Suppose that $(M, d)$ is compact and that $\rho$ is another metric on $M$.
    Then $d, \rho$ are equivalent iff they are uniformly equivalent.
\end{proposition}

\begin{theorem}{8.20}
    Let $(V, \lvert \lvert \cdot \rvert \rvert), (V, \lvert \lvert \lvert \cdot \rvert \rvert \rvert)$ be normed vector spaces, and let $T: V \rightarrow W$ be a linear map.
    Then the following are equivalent:
    
    (i) $T$ is Lipschitz

    (ii) $T$ is uniformly continuous

    (iii) $T$ is continuous everywhere

    (iv) $T$ is continuous at $0 \in V$

    (v) there is a constant $C$ such that $\lvert \lvert \lvert T(x) \rvert \rvert \rvert \leq C \lvert \lvert x \rvert \rvert$ for all $x \in V$
\end{theorem}